% A6 Vector guide. Build with: xelatex Guide_A6.tex
\documentclass[11pt,a4paper]{article}

\usepackage[margin=2.2cm]{geometry}
\usepackage{fontspec}
\usepackage{xeCJK}
\setmainfont{TeX Gyre Heros}
\setmonofont{DejaVu Sans Mono}[Scale=0.88]
\setCJKmainfont{WenQuanYi Zen Hei}[AutoFakeBold=2.5]
\setCJKmonofont{WenQuanYi Zen Hei}
\xeCJKsetup{CJKecglue={\hskip 0.2em plus 0.1em}}

\usepackage{xcolor}
\usepackage{titlesec}
\usepackage{booktabs}
\usepackage{tabularx}
\usepackage{array}
\usepackage{enumitem}
\usepackage{listings}
\usepackage[most]{tcolorbox}
\usepackage[hidelinks]{hyperref}

\definecolor{accent}{HTML}{1F5FAD}
\definecolor{codebg}{HTML}{F5F7FA}
\definecolor{donebg}{HTML}{EEF6EE}
\definecolor{donefr}{HTML}{3C8D40}
\definecolor{warnbg}{HTML}{FFF6E5}
\definecolor{warnfr}{HTML}{D08A00}

\titleformat{\section}{\Large\bfseries\color{accent}}{}{0pt}{}
\titleformat{\subsection}{\large\bfseries}{}{0pt}{}
\titlespacing*{\section}{0pt}{2.2ex plus .5ex}{1ex}
\titlespacing*{\subsection}{0pt}{1.8ex plus .4ex}{0.8ex}
\setlength{\parindent}{0pt}
\setlength{\parskip}{0.7ex}
\linespread{1.15}
\setlist{nosep, leftmargin=1.6em, topsep=0.3ex}
\renewcommand{\arraystretch}{1.25}

\lstset{
  language=C++, basicstyle=\ttfamily\small, columns=fullflexible, keepspaces=true,
  backgroundcolor=\color{codebg}, frame=single, rulecolor=\color{codebg},
  framesep=4pt, xleftmargin=4pt, xrightmargin=4pt, aboveskip=0.8ex, belowskip=0.8ex,
  keywordstyle=\color{accent}\bfseries, commentstyle=\color{gray}, stringstyle=\color{donefr},
  showstringspaces=false, morekeywords={nullptr}
}

\newcommand{\code}[1]{\texttt{#1}}
\newcommand{\idea}{\textbf{思路：}}
\newcommand{\todo}{\textbf{你来写：}}

\newtcolorbox{done}{enhanced, before skip=2.2ex, colback=donebg, colframe=donefr, boxrule=0.6pt, arc=2pt,
  left=6pt, right=6pt, top=8pt, bottom=4pt, fontupper=\small, before upper=\raggedright,
  title={完成标志}, coltitle=white, fonttitle=\small\bfseries,
  attach boxed title to top left={xshift=6pt, yshift=-2mm},
  boxed title style={colback=donefr, colframe=donefr, arc=1pt, boxrule=0pt}}
\newtcolorbox{tip}[1][小知识]{enhanced, before skip=2.2ex, colback=warnbg, colframe=warnfr, boxrule=0.6pt, arc=2pt,
  left=6pt, right=6pt, top=8pt, bottom=4pt, fontupper=\small, before upper=\raggedright,
  title={#1}, coltitle=white, fonttitle=\small\bfseries,
  attach boxed title to top left={xshift=6pt, yshift=-2mm},
  boxed title style={colback=warnfr, colframe=warnfr, arc=1pt, boxrule=0pt}}

\begin{document}
\sloppy
\setlength{\emergencystretch}{3em}

\begin{center}
  {\LARGE\bfseries A6 Vector：动手指南}\\[0.8ex]
  {\normalsize Assignment 006 · 类模板 Vector}
\end{center}

\textbf{目标：}按题目给出的声明，实现一个自己的动态数组类模板 \code{Vector<T>}，并写一个 \code{main} 测试它的全部功能。

本文是一条推荐路线，不是唯一写法。每一步写完都要编译、跑一下“完成标志”里的例子，通过了再做下一步。这次评分第三项是 \textbf{Test coverage}（测试覆盖），所以测试和代码同样重要。

\section{一、题目要求对照}

题目已经给出了类的声明，你要写的是每个成员函数的函数体。先弄清每个成员做什么：

\begin{center}
\small
\begin{tabularx}{\linewidth}{@{}>{\raggedright\arraybackslash}p{4.4cm}>{\raggedright\arraybackslash}X@{}}
\toprule
成员 & 要做的事 \\
\midrule
\code{Vector()} & 空向量：不分配内存，大小和容量都是 0 \\
\code{Vector(int size)} & 准备好 \code{size} 个元素（含义见第二节约定 2） \\
\code{Vector(const Vector\& r)} & 拷贝构造：\textbf{深拷贝}，复制一份自己的数组 \\
\code{\textasciitilde Vector()} & 析构：释放数组 \\
\code{operator[](int)} & 按下标访问，\textbf{不}检查越界 \\
\code{at(int)} & 按下标访问，越界时抛 \code{std::out\_of\_range} \\
\code{size()}、\code{empty()} & 返回元素个数；是否为空 \\
\code{push\_back(const T\&)} & 在末尾加一个元素，空间不够先扩容 \\
\code{clear()} & 清空内容 \\
\code{inflate()}（私有） & 扩容：申请更大的数组，搬运元素，释放旧数组 \\
\bottomrule
\end{tabularx}
\end{center}

三个数据成员的关系是理解整个作业的关键：

\begin{center}
\small
\begin{tabular}{@{}ll@{}}
\toprule
成员 & 含义 \\
\midrule
\code{m\_pElements} & 指向 \code{new} 出来的数组 \\
\code{m\_nCapacity} & 数组一共有几个格子（容量） \\
\code{m\_nSize} & 其中前几个格子放了真正的元素（大小），始终 $0 \le$ \code{m\_nSize} $\le$ \code{m\_nCapacity} \\
\bottomrule
\end{tabular}
\end{center}

\begin{tip}[一句话理解 Vector]
Vector 就是“一块会自己搬家的数组”：格子用完了，就换一块两倍大的新数组，把旧元素搬过去，再把旧数组还给系统。
\end{tip}

\section{二、先定几条约定}

\subsection{1. 模板的代码全部写在头文件里}

类模板不是真正的类，而是“生成类的模具”。编译器只有在看到 \code{Vector<int>} 这样的用法\textbf{并且}同时看得到函数体时，才会生成对应的代码。

如果像以前一样把函数体放进 \code{Vector.cpp}，编译 \code{Vector.cpp} 时没有人用到 \code{Vector<int>}，什么都不会生成；编译 \code{main.cpp} 时又看不到函数体。结果就是链接错误（已实测）：

\begin{lstlisting}[language={}]
undefined reference to `Vector<int>::push_back(int const&)'
\end{lstlisting}

所以推荐：\code{Vector.h} 里先放题目给的类声明，紧接着在\textbf{同一个文件}里写所有成员函数的定义。

\subsection{2. Vector(int size) 的含义}

题目注释是 “creates a vector for holding 'size' elements”，有两种理解：

\begin{center}
\small
\begin{tabularx}{\linewidth}{@{}lXX@{}}
\toprule
 & 理解 A（推荐） & 理解 B \\
\midrule
做法 & 像 \code{std::vector<int> v(5)}：里面已经有 \code{size} 个元素，值为 0 或默认值 & 只预留 \code{size} 个格子，里面还没有元素 \\
\code{m\_nSize} & \code{size} & 0 \\
\code{m\_nCapacity} & \code{size} & \code{size} \\
构造后 \code{v[2] = 1} & 合法 & 越界 \\
\bottomrule
\end{tabularx}
\end{center}

推荐理解 A，和标准库行为一致，测试也更直观。选哪种都行，但要在 README 里写清楚。\code{size} 为负数时，抛 \code{std::invalid\_argument}。

\subsection{3. 扩容策略}

容量为 0 时扩成 1，否则扩成原来的 2 倍（题目注释里的 \code{2*m\_nCapacity}）。翻倍让 \code{push\_back} 平均下来很快：放 1000 个元素只需要扩容 11 次。

\subsection{4. 补一个赋值运算符}

题目的声明里有拷贝构造和析构，却\textbf{没有}拷贝赋值运算符 \code{operator=}。编译器会自动生成一个“浅拷贝”版本，只复制指针。于是：

\begin{lstlisting}
Vector<int> a, b;
b.push_back(1);
a = b;          // a and b now point to the SAME array
                // at the end of main both destructors delete[] it
\end{lstlisting}

用地址检查工具（AddressSanitizer）实测，程序结束时报 \code{attempting double-free}（同一块内存被释放两次）。\code{-Wextra} 也会对 \code{a = b} 报 \code{-Wdeprecated-copy} 警告。

这就是 C++ 的“三法则”（Rule of Three）：\textbf{自己写了析构、拷贝构造、拷贝赋值中的任何一个，通常三个都要写。}推荐在 \code{public} 里补一行声明 \code{Vector\& operator=(const Vector\& r);}，并在 README 里说明“为避免浅拷贝，在题目声明之外补充了赋值运算符”。

\section{三、代码怎么分}

\begin{center}
\small
\begin{tabular}{@{}ll@{}}
\toprule
文件 & 内容 \\
\midrule
\code{Vector.h} & 题目的类声明（加上 \code{operator=}），以及所有成员函数的定义 \\
\code{main.cpp} & 测试：逐项调用并打印 PASS / FAIL \\
\bottomrule
\end{tabular}
\end{center}

\textbf{类外定义模板成员函数的写法}是这次的新语法。每个函数前面都要加 \code{template <class T>}，类名写成 \code{Vector<T>::}。示范默认构造和 \code{size()}：

\begin{lstlisting}
template <class T>
Vector<T>::Vector() : m_pElements(nullptr), m_nSize(0), m_nCapacity(0) {}

template <class T>
int Vector<T>::size() const {
    return m_nSize;
}
\end{lstlisting}

返回类型里也要写全，例如赋值运算符的定义是 \code{template <class T> Vector<T>\& Vector<T>::operator=(const Vector\& r)}。参数列表里的 \code{Vector} 可以省略 \code{<T>}，因为已经进入了类的作用域。

\section{四、按五步实现}

\subsection{第 1 步：构造与析构}

\idea 先让对象能正确地“出生”和“死亡”，不泄漏内存。

\todo
\begin{itemize}
  \item 默认构造：照上面的示范。
  \item \code{Vector(int size)}：负数抛异常；大于 0 时用 \code{new T[size]()} 申请数组（末尾的 \code{()} 让 \code{int} 等基本类型初始化为 0）；大小和容量都设为 \code{size}。
  \item 析构：\code{delete[] m\_pElements;}。
\end{itemize}

\begin{tip}
\begin{itemize}[leftmargin=1.2em]
  \item \code{new T[n]} 必须配 \code{delete[]}，不能写成 \code{delete}。
  \item \code{delete[] nullptr} 是安全的，什么也不做，所以空向量析构不用特判。
  \item \code{new T[n]} 要求 \code{T} 有默认构造函数。本作业的测试类型（\code{int}、\code{std::string}）都满足。
\end{itemize}
\end{tip}

\begin{done}
\code{Vector<int> e;} 的 \code{size()} 为 0；\code{Vector<int> z(3);} 的 \code{size()} 为 3，三个元素都是 0；\code{Vector<int> bad(-1);} 抛出 \code{std::invalid\_argument}。
\end{done}

\subsection{第 2 步：size、empty、[] 和 at}

\todo
\begin{itemize}
  \item \code{size()} 返回 \code{m\_nSize}；\code{empty()} 返回 \code{m\_nSize == 0}。
  \item \code{operator[]}：直接 \code{return m\_pElements[index];}，题目明确要求不检查。
  \item \code{at}：\code{index < 0 || index >= m\_nSize} 时 \code{throw std::out\_of\_range("...")}（头文件 \code{<stdexcept>}），否则同上。
\end{itemize}

注意 \code{at} 比较的是 \code{m\_nSize} 而不是 \code{m\_nCapacity}：容量里多出来的格子还没有放元素，不能访问。

\begin{tip}[为什么返回 T\&]
返回引用，调用者拿到的就是数组里的那个元素本身，所以 \code{v[0] = 5;} 能修改向量里的值。如果返回 \code{T}，改的只是一份副本。
\end{tip}

\begin{done}
\code{Vector<int> z(3); z[1] = 7;} 之后 \code{z.at(1)} 为 7；\code{z.at(3)} 和 \code{z.at(-1)} 都抛出 \code{std::out\_of\_range}。
\end{done}

\subsection{第 3 步：inflate 和 push\_back}

\idea \code{push\_back} 只做两件事：满了就调用 \code{inflate()}，然后把元素放到下标 \code{m\_nSize} 的位置，大小加一。

\todo \code{inflate()} 按顺序做：
\begin{enumerate}
  \item 算出新容量（约定 3）；
  \item \code{new} 一个新数组；
  \item 用循环把前 \code{m\_nSize} 个元素复制过去；
  \item \code{delete[]} 旧数组；
  \item 让 \code{m\_pElements} 指向新数组，更新 \code{m\_nCapacity}。
\end{enumerate}

\begin{tip}[一个隐蔽的 bug：push\_back 自己的元素]
\code{v.push\_back(v[0]);} 时，参数 \code{x} 是对 \code{v[0]} 的引用，指向\textbf{旧数组}。如果这时恰好要扩容，\code{inflate()} 已经把旧数组 \code{delete[]} 了，再读 \code{x} 就是读已释放的内存。实测 AddressSanitizer 报 \code{heap-use-after-free}。

修法很简单：在 \code{push\_back} 一开始先复制一份 \code{T value = x;}，再扩容，最后放入 \code{value}。
\end{tip}

\begin{done}
循环 \code{push\_back} 0 到 999，\code{size()} 为 1000，每个 \code{v[i] == i}；只有一个元素的向量连续 \code{push\_back(one[0])} 20 次，结果全是原来的值，地址检查不报错。
\end{done}

\subsection{第 4 步：clear}

\todo 把 \code{m\_nSize} 设为 0 即可，保留已经申请的数组，之后 \code{push\_back} 可以直接复用，不用重新扩容。这也是 \code{std::vector::clear} 的做法。

\begin{done}
\code{clear()} 后 \code{empty()} 为真；再 \code{push\_back(3)}，\code{size()} 为 1，\code{v[0]} 为 3。
\end{done}

\subsection{第 5 步：拷贝构造和赋值运算符}

\idea 每个 Vector 都要有\textbf{自己的}数组。拷贝时复制元素，而不是复制指针。

\todo
\begin{itemize}
  \item 拷贝构造：按 \code{r.m\_nSize} 申请新数组（为 0 时保持 \code{nullptr}），循环复制元素，设置大小和容量。
  \item 赋值运算符：先判断 \code{this != \&r}（防止 \code{a = a} 时先把自己的数组删掉），再做和拷贝构造类似的事，记得先 \code{delete[]} 自己原来的数组，最后 \code{return *this;}。
\end{itemize}

\begin{tip}[进阶写法：copy-and-swap]
赋值运算符也可以借用拷贝构造：先 \code{Vector copy(r);}，再用 \code{std::swap}（头文件 \code{<utility>}）交换 \code{*this} 和 \code{copy} 的三个成员。函数结束时 \code{copy} 析构，顺便释放了旧数组。不用重复写复制循环，即使 \code{new} 失败，原对象也不会被破坏。
\end{tip}

\begin{done}
\code{Vector<int> c(v); c[0] = -7;} 之后 \code{v[0]} 不变；\code{a = v; a[1] = -1;} 之后 \code{v[1]} 不变；\code{a = a;} 后内容不变；拷贝一个空向量得到空向量。
\end{done}

\section{五、测试}

这次测试覆盖是评分项，要让每个成员函数都被调用到，并覆盖边界情况。沿用 A5 的 \code{check} 小函数：

\begin{lstlisting}
int failures = 0;

void check(bool ok, const std::string& name) {
    if (!ok) ++failures;
    std::cout << (ok ? "PASS  " : "FAIL  ") << name << '\n';
}
\end{lstlisting}

\begin{center}
\small
\begin{tabularx}{\linewidth}{@{}lX@{}}
\toprule
类别 & 用例 \\
\midrule
构造 & 默认构造；\code{Vector(3)} 的大小和初值；\code{Vector(-1)} 抛异常 \\
访问 & \code{[]} 读写；\code{at} 读写；\code{at(-1)}、\code{at(size())} 抛 \code{std::out\_of\_range} \\
增长 & 连续 \code{push\_back} 1000 个并逐个核对；\code{push\_back} 自己的元素 \\
清空 & \code{clear} 后为空；清空后再 \code{push\_back} \\
拷贝 & 拷贝构造和赋值后修改副本，原对象不变；自我赋值；拷贝空向量 \\
其他类型 & \code{Vector<std::string>}（会分配内存的类型）；\code{Vector<Vector<int>>}（检验深拷贝是否正确） \\
\bottomrule
\end{tabularx}
\end{center}

异常用例用 \code{try / catch}：没抛异常算 FAIL，抛了正确的类型算 PASS。

内存错误（泄漏、重复释放、访问已释放内存）往往不会让程序立刻崩溃，肉眼看不出来。编译时加上地址检查，有问题会直接报错：

\begin{lstlisting}[language=bash]
g++ -std=c++17 -Wall -Wextra -g -fsanitize=address,undefined main.cpp -o vector_test
./vector_test
\end{lstlisting}

模板全在头文件里，所以只需要编译 \code{main.cpp}。全部 PASS、编译没有警告、地址检查没有报错，才算完成。

\begin{tip}[Windows 用户]
MinGW 的 g++ 不支持 \code{-fsanitize=address}。可以在 WSL 里编译运行，或者去掉这个选项，只看 PASS / FAIL。
\end{tip}

\section{六、最后提交什么}

题目要求提交包含全部源文件的 \code{.zip}。按课程统一规范：

\begin{itemize}
  \item \code{code/}：\code{Vector.h}、\code{main.cpp}；不放可执行文件。
  \item \code{README.md}：环境、编译和运行命令、你的约定：\code{Vector(int)} 的含义、扩容策略、补充的 \code{operator=}、哪些情况抛异常。
  \item \code{report/}：按课程命名的 PDF，包含环境、运行方式和真实的测试结果。
\end{itemize}

把它们放进 \code{Assignment06\_<学号>\_<姓名拼音>/}，打包成 \code{.zip}。

\bigskip
\begin{center}
\textbf{现在只做第 1 步：写出默认构造、Vector(int) 和析构，验证大小和初值。通过后再开始第 2 步。}
\end{center}

\end{document}
